\section{Protocolos de comunicación}

\subsection{I2C}
La comunicación principal entre los ATmega328P utiliza I2C. El Maestro configura el bus a 100 kHz y los Slaves responden en las direcciones 0x11 y 0x12. El VL53L0X, denominado Slave 3, comparte las mismas líneas SDA y SCL con dirección 0x29. SDA corresponde a PC4/A4 y SCL a PC5/A5 en los tres Arduino Nano.

Sobre I2C se implementó un protocolo propio de solicitudes y respuestas. El Maestro puede verificar la disponibilidad de los nodos, solicitar sensores y enviar órdenes a los actuadores. Los periféricos decodifican las solicitudes fuera de la rutina de interrupción TWI y publican respuestas que incluyen estados como operación correcta, solicitud aceptada, dispositivo no soportado, dispositivo desconocido o ausencia de una muestra válida. Esta estructura permite separar el transporte I2C de las acciones específicas del Car Wash.

El Maestro también comprueba periódicamente la comunicación con ambos Slaves. Durante la operación normal realiza una verificación aproximadamente cada 500 ms y entra a un estado de error si se pierde uno de los enlaces.

\subsection{UART}
El Maestro configura el USART del ATmega328P a 9600 baudios con formato de 8 bits de datos, sin paridad y un bit de parada. PD0/RX y PD1/TX quedan destinados a esta interfaz. La transmisión utiliza un buffer y la interrupción de registro de datos vacío, evitando que la impresión de mensajes bloquee directamente la lógica principal.

UART cumple dos funciones. La primera es diagnóstico y supervisión: el Maestro informa transiciones de estado, distancias, errores y permite solicitar reportes o ejecutar un escaneo del bus I2C. La segunda es el enlace con el ESP32. El Maestro transmite por PD1/D1 tramas ASCII de telemetría a 9600 baudios, mientras que las órdenes provenientes del ESP32 ingresan por PB4/D12 mediante una UART por software. En el ESP32 se utiliza UART2, con GPIO16 como RX y GPIO17 como TX.

Las tramas de telemetría siguen el formato \texttt{@CW,T,CLAVE,VALOR}, mientras que las órdenes remotas utilizan \texttt{@CW,C,DISPOSITIVO,VALOR}. Esta separación permite integrar la supervisión de Adafruit IO sin reemplazar el protocolo I2C utilizado entre los ATmega328P.
